\documentclass[aspectratio=169,professionalfonts]{beamer}
\usepackage[T1]{fontenc}
\usepackage[utf8]{inputenc}
\usepackage{lmodern}
\usepackage{amsmath,amssymb}
\usepackage{stmaryrd}
\usepackage{xcolor}
\usepackage{pifont}
\usepackage{tikz}
\usetheme{default}
\setbeamertemplate{navigation symbols}{}
\setbeamertemplate{headline}{}
\setbeamersize{text margin left=0.95cm,text margin right=0.95cm}

% High-contrast chalkboard palette
\definecolor{Board}{HTML}{052B22}
\definecolor{BoardDeep}{HTML}{031E18}
\definecolor{Chalk}{HTML}{F8F2DF}
\definecolor{SoftChalk}{HTML}{EFE6C9}
\definecolor{BrightMint}{HTML}{D7FFD6}
\definecolor{Sage}{HTML}{A8D5B3}
\definecolor{GoldChalk}{HTML}{E3C173}
\definecolor{Dust}{HTML}{C7D4C1}
\definecolor{AlertGold}{HTML}{FFD06A}
\definecolor{Auburn}{HTML}{A86B38}
\definecolor{ChalkBlue}{HTML}{A9CEFF}
\definecolor{ChalkPeach}{HTML}{FFC49B}
\definecolor{ChalkRose}{HTML}{F0A8C0}
\definecolor{ChalkLilac}{HTML}{C9B8F5}

\setbeamercolor{background canvas}{bg=Board}
\setbeamercolor{normal text}{fg=Chalk,bg=Board}
\setbeamercolor{structure}{fg=BrightMint}
\setbeamercolor{frametitle}{fg=Chalk,bg=Board}
\setbeamercolor{title}{fg=Chalk}
\setbeamercolor{subtitle}{fg=Sage}
\setbeamercolor{author}{fg=SoftChalk}
\setbeamercolor{date}{fg=SoftChalk}
\setbeamercolor{math text}{fg=BrightMint}
\setbeamercolor{math text displayed}{fg=BrightMint}
\setbeamercolor{itemize item}{fg=GoldChalk}
\setbeamercolor{itemize subitem}{fg=Sage}
\setbeamercolor{aibox}{fg=Chalk,bg=BoardDeep}
\setbeamercolor{userbox}{fg=Chalk,bg=BoardDeep}
\setbeamercolor{warnbox}{fg=Chalk,bg=BoardDeep}

\setbeamerfont{title}{series=\bfseries,size=\fontsize{30}{36}\selectfont}
\setbeamerfont{author}{size=\fontsize{16}{20}\selectfont}
\setbeamerfont{frametitle}{series=\bfseries,size=\fontsize{20}{24}\selectfont}
\setbeamerfont{normal text}{size=\normalsize}

\setbeamertemplate{frametitle}{%
  \vspace{0.28cm}%
  \insertframetitle\par%
  \vspace{0.08cm}%
  {\color{GoldChalk}\rule{0.22\paperwidth}{0.75pt}}%
  \vspace{0.10cm}%
}

\setbeamertemplate{footline}{%
  \leavevmode%
  \hbox{%
    \begin{beamercolorbox}[wd=.80\paperwidth,ht=2.35ex,dp=0.8ex,leftskip=1.2em]{normal text}%
      \color{Sage}\scriptsize Getting leverage on normativity
    \end{beamercolorbox}%
    \begin{beamercolorbox}[wd=.20\paperwidth,ht=2.35ex,dp=0.8ex,rightskip=1.2em plus1fil]{normal text}%
      \hfill\color{Sage}\scriptsize \insertframenumber/\inserttotalframenumber
    \end{beamercolorbox}%
  }%
}

\newenvironment{aibox}{%
  \begin{beamercolorbox}[rounded=true,shadow=false,wd=\textwidth,sep=0.65ex]{aibox}%
  \fontsize{10.4}{12.4}\selectfont\color{Chalk}\textbf{\color{AlertGold}AI-draft scaffold}\par\vspace{0.15em}%
}{\end{beamercolorbox}}

\newenvironment{userbox}{%
  \begin{beamercolorbox}[rounded=true,shadow=false,wd=\textwidth,sep=0.65ex]{userbox}%
  \fontsize{10.4}{12.4}\selectfont\color{Chalk}\textbf{\color{BrightMint}From Anson draft / reviewed language}\par\vspace{0.15em}%
}{\end{beamercolorbox}}

\newenvironment{planbox}{%
  \begin{beamercolorbox}[rounded=true,shadow=false,wd=\textwidth,sep=0.45ex]{aibox}%
  \fontsize{8}{9.8}\selectfont\color{Dust}\textbf{\color{Auburn}AI scaffold --- reveal plan}\quad}{\end{beamercolorbox}}

\newenvironment{ansonbox}[1]{%
  \begin{beamercolorbox}[rounded=true,shadow=false,wd=\textwidth,sep=0.65ex]{userbox}%
  \fontsize{13.2}{16.2}\selectfont\color{Chalk}\textbf{\color{SoftChalk}#1}\par\vspace{0.2em}%
}{\end{beamercolorbox}}
\newenvironment{plainbox}{%
  \begin{beamercolorbox}[rounded=true,shadow=false,wd=\textwidth,sep=0.65ex]{userbox}%
  \fontsize{13.2}{16.2}\selectfont\color{Chalk}%
}{\end{beamercolorbox}}

\makeatletter
\newenvironment{outlinebox}[1]{%
  \noindent\setlength{\fboxsep}{5pt}\setlength{\fboxrule}{1pt}%
  \begin{lrbox}{\@tempboxa}\begin{minipage}{\dimexpr\textwidth-2\fboxsep-2\fboxrule\relax}%
  \fontsize{13.2}{16.2}\selectfont\color{Chalk}\textbf{\color{SoftChalk}#1}\par\vspace{0.25em}%
}{\end{minipage}\end{lrbox}\fcolorbox{SoftChalk}{Board}{\usebox{\@tempboxa}}}
\makeatother
\makeatletter
\newenvironment{outlineplain}{%
  \noindent\setlength{\fboxsep}{6pt}\setlength{\fboxrule}{1pt}%
  \begin{lrbox}{\@tempboxa}\begin{minipage}{\dimexpr\textwidth-2\fboxsep-2\fboxrule\relax}%
  \fontsize{13.2}{16.2}\selectfont\color{Chalk}%
}{\end{minipage}\end{lrbox}\fcolorbox{SoftChalk}{Board}{\usebox{\@tempboxa}}}
\makeatother
\newcommand{\mintline}[1]{{\color{BrightMint}\bfseries #1}}
\newlength{\levboxht}\setlength{\levboxht}{106pt}

\newcommand{\keyline}[1]{{\color{GoldChalk}\bfseries #1}}
\newcommand{\drafttag}{\textbf{\color{AlertGold}[AI-DRAFT]}}
% Corner badge marking a frame whose on-slide text is currently LLM-written.
% Delete the \llmtag line from a frame once the wording is Anson's.
\newcommand{\llmtag}{%
\begin{tikzpicture}[remember picture,overlay]
\node[anchor=north east,xshift=-0.50cm,yshift=-0.34cm] at (current page.north east)
{\setlength{\fboxsep}{2.8pt}\setlength{\fboxrule}{0.9pt}%
\fcolorbox{AlertGold}{BoardDeep}{\scriptsize\bfseries\color{AlertGold}LLM-drafted text}};
\end{tikzpicture}%
\par\vspace*{-\baselineskip}}
\newcommand{\eqtag}[1]{{\color{Sage}\text{#1}}}
% Statics-recap pipeline connectives. Set per side so the *rendered ink*
% gaps come out equal (glyph side bearings differ left vs right).
\newcommand{\pipegapL}{21.76mu}
\newcommand{\pipegapR}{5.29mu}
\newcommand{\withgapL}{21.76mu}
\newcommand{\withgapR}{7.15mu}

\title{Getting Leverage on Normativity}
\author{Anson Berns}
\date{}

\begin{document}

\begin{frame}[plain,noframenumbering]
  \vspace{1.55cm}
  {\usebeamerfont{title}\color{Chalk}Getting Leverage on Normativity\par}
  \vspace{0.28cm}
  {\color{GoldChalk}\rule{0.38\paperwidth}{1.0pt}\par}
  \vspace{0.58cm}
  {\usebeamerfont{author}\color{SoftChalk}Anson Berns\par}
  \vfill
  {\color{Sage}\fontsize{13.5}{16}\selectfont Draft 8/26}
\end{frame}

\begin{frame}[t,shrink=2]{Motivation: Learning Normativity}
\begin{ansonbox}{Learning Normativity desiderata\hfill{\normalfont\color{Dust}\small Demski 2020}}
\begin{itemize}\setlength\itemsep{0.4em}
  \item Whole-process, all-levels learning
  \item The meaning of any feedback can later be reinterpreted
  \item No final, immutable fact of the matter to discover
\end{itemize}
\end{ansonbox}
\vspace{0.8em}
\onslide<2->{{\large \textbf{\color{SoftChalk}Goal:} build a solution to these desiderata using logical induction as a key ingredient\par}}
\end{frame}

\begin{frame}[t]{The law}
\begin{ansonbox}{A paradigm domain for Learning Normativity}
\begin{itemize}\setlength\itemsep{0.4em}
  \item Judicial decisions rule on all levels of legal reasoning
  \item Rulings can always be appealed or overturned
  \item Even the highest court might be wrong about a particular verdict
\end{itemize}
\end{ansonbox}
\vspace{0.65em}
\onslide<2->{%
{\large It's not enough to predict human judges, we want to learn the underlying norm.\par}
\vspace{0.45em}
{\large\mbox{How should AI put a credence on \keyline{``John Smith should be found guilty''}?}\par}}
\end{frame}

\begin{frame}[t,shrink=2]{A middle way}
\vspace{0.45cm}
{\large Even though there's \keyline{no ``final answer''}, our judgments are \keyline{not arbitrary}.\par}
\vspace{0.5cm}
\begin{ansonbox}{Example}
Reasoned disagreement is always possible about a legal case, but some judges are better than others
\end{ansonbox}
\vspace{0.75cm}
\onslide<2->{{\large \mintline{Leverage}: how are normative beliefs \keyline{constrained}, without being pinned to a single number?\par}}
\end{frame}

\begin{frame}[t]{Two kinds of leverage}
\begin{columns}[T,onlytextwidth]
\begin{column}{0.478\textwidth}
\begin{ansonbox}{Coherence}
\begin{minipage}[t][\levboxht][t]{\dimexpr\linewidth-6.5pt\relax}
How do normative claims bear on each other?
\onslide<2->{%
\begin{itemize}\setlength\itemsep{0.25em}
  \item Similar cases get similar verdicts
  \item Doctrine and precedent
  \item Internal logic of the law
\end{itemize}
\smallskip
\onslide<3->{\mbox{\keyline{More than logical consistency}}}
}
\end{minipage}
\end{ansonbox}
\end{column}
\begin{column}{0.478\textwidth}
\begin{ansonbox}{Groundedness}
\begin{minipage}[t][\levboxht][t]{\dimexpr\linewidth-6.5pt\relax}
How do empirical claims bear on normative claims?
\onslide<4->{%
\begin{itemize}\setlength\itemsep{0.25em}
  \item Physical evidence
  \item Testimony
  \item History and context of the legal code
\end{itemize}
}
\end{minipage}
\end{ansonbox}
\end{column}
\end{columns}
\vspace{0.95em}
\onslide<5->{{\large Particular coherence and groundedness claims are themselves up for debate\par}}
\end{frame}

\begin{frame}[t]{What logical induction is missing}
\only<1-3>{%
\vspace{0.55cm}
{\fontsize{13.2}{16.8}\selectfont
\begin{itemize}\setlength\itemsep{1.1em}
  \item<1-> All the normative action happens with \keyline{never-settling} sentences, but settlement is what constrains prices in LI
  \item<2-> The only kind of external \keyline{feedback} available to a vanilla logical inductor is immutable certainty from deduction
  \item<3-> We are lacking a representation of \keyline{reasons}, the network of non-logical relationships between sentences
\end{itemize}
}
}%
\only<4>{%
\vspace{1.5cm}
\begin{center}
{\fontsize{16}{22}\selectfont
As a solution to the Learning Normativity desiderata, logical induction is missing an account of \keyline{leverage}.
\par}
\end{center}
}%
\end{frame}

\begin{frame}[t]{An afoundationalist picture}
\only<1>{%
\vspace{0.35cm}
\begin{outlinebox}{Foundationalist accounts}
\begin{itemize}\setlength\itemsep{0.5em}
  \item \textbf{\color{SoftChalk}Metaphor:} A layer of immutable bedrock underpinning every story of a tall building
  \item Based on \keyline{self-justifying certainty}, deepest principles that can't be questioned
\end{itemize}
\end{outlinebox}
}%
\only<2->{%
\vspace{0.25cm}
{\fontsize{13.2}{16.8}\selectfont
\begin{itemize}\setlength\itemsep{0.7em}
  \item<2-> We need a new picture that's neither foundationalist nor flimsy
  \item<3-> \textbf{\color{SoftChalk}Afoundationalist metaphor:} An impermanent seed that constrains the growth of the eventual tree without predetermining it
  \item<3-> Based on \keyline{revisable pre-licensed entitlements}, an initial induction into the normative domain that can be questioned after learning
  \item<4-> Defeasible bridge principles that let us go beyond coherentism
\end{itemize}
}
}%
\end{frame}

\begin{frame}[t]{A bounded reasoner}
\vspace*{0.10cm}
{\fontsize{13.2}{16.8}\selectfont
\begin{itemize}\setlength\itemsep{0.5em}
  \item<1-> Take a language $\mathcal L$, a set of logical, empirical, and normative sentences
  \item<1-> For now, fix a time $t$ and the finite subset $\mathcal L_t\subset\mathcal L$ of currently relevant sentences
  \item<2-> Let $\Omega_t$ be the possible worlds consistent with the settled record $D_t$ for the currently relevant sentences
\end{itemize}
}
\vspace{0.05cm}
\onslide<3->{%
\begin{ansonbox}{Credence states}
\smallskip
\centering
$\displaystyle \mu\in\Delta(\Omega_t)$
\qquad\quad
$\displaystyle \mu(\varphi)\;=\!\!\sum_{\substack{\omega\in\Omega_t\\ \text{with }\omega\models\varphi}}\!\!\mu(\omega)$
\par
\end{ansonbox}
}
\end{frame}

\begin{frame}[t]{Warrants}
\vspace{0.35cm}
{\fontsize{13.2}{16.8}\selectfont
\begin{itemize}\setlength\itemsep{0.6em}
  \item<1-> \textbf{\color{SoftChalk}New primitive:} a \keyline{warrant} $\rho:\ \alpha(n)\rightsquigarrow\beta(n)$
  \begin{itemize}\setlength\itemsep{0.45em}
    \item<1-> ``$\alpha$ sentences provide support for $\beta$ sentences because of defeasible reason $\rho$''
    \item<1-> One schema that applies to multiple cases, indexed by $n$
  \end{itemize}
  \vspace{0.55em}
  \item<2-> Each warrant gets an \keyline{applicability sentence} $\mathsf{App}_\rho(n)$ in the language $\mathcal L$
  \begin{itemize}\setlength\itemsep{0.45em}
    \item<2-> ``Warrant $\rho$ really applies in case $n$''
  \end{itemize}
  \vspace{0.55em}
  \item<3-> Warrants themselves are strengthless, numbers are going to enter contextually (soon)
\end{itemize}
}
\end{frame}

\begin{frame}[t]{Undercutting and rebuttal}
\vspace{0.3cm}
{\fontsize{13.2}{16.8}\selectfont
This Horty-inspired setup distinguishes two kinds of \keyline{conflict} between reasons
\par}
\vspace{0.35cm}
\begin{columns}[t]
\begin{column}{0.478\textwidth}
\begin{ansonbox}{Rebuttal}
\centering
$\rho:\ {\color{ChalkPeach}\alpha(n)}\rightsquigarrow{\color{ChalkBlue}\beta(n)}$
\par\medskip
$\rho':\ {\color{ChalkRose}\gamma(n)}\rightsquigarrow\lnot{\color{ChalkBlue}\beta(n)}$
\par
\end{ansonbox}
\onslide<2->{%
\vspace{0.6cm}
{\fontsize{12}{15}\selectfont\centering
\begin{tabular}{@{}c@{}}
{\color{ChalkPeach}$\alpha$: ``Alice says she saw John do it''}\\[0.5em]
{\color{ChalkBlue}$\beta$: ``John should be found guilty''}\\[0.5em]
{\color{ChalkRose}$\gamma$: ``Bob says John was with him''}
\end{tabular}\par}
}
\end{column}
\begin{column}{0.478\textwidth}
\begin{ansonbox}{Undercutting}
\centering
$\rho:\ {\color{ChalkPeach}\alpha(n)}\rightsquigarrow{\color{ChalkBlue}\beta(n)}$
\par\medskip
$\rho'':\ {\color{ChalkLilac}\delta(n)}\rightsquigarrow\lnot\mathsf{App}_\rho(n)$
\par
\end{ansonbox}
\onslide<2->{%
\vspace{0.6cm}
{\fontsize{12}{15}\selectfont\centering
\begin{tabular}{@{}c@{}}
{\color{ChalkPeach}$\alpha$: ``Alice says she saw John do it''}\\[0.5em]
{\color{ChalkBlue}$\beta$: ``John should be found guilty''}\\[0.5em]
{\color{ChalkLilac}$\delta$: ``Alice was drunk''}
\end{tabular}\par}
}
\end{column}
\end{columns}
\end{frame}

\begin{frame}[t]{Endorsements}
\vspace{-0.4cm}
{\fontsize{13.2}{16.8}\selectfont
Now, let's add numbers:
\par}
\vspace{0.15cm}
{\fontsize{13.2}{16.8}\selectfont
\begin{itemize}\setlength\itemsep{0.6em}
  \item<1-> An \keyline{endorsement} quantifies a warrant in a particular context
  \item<2-> An inequality on credences like $\mu(\varphi\wedge\psi)\ge c\cdot\mu(\varphi)$, with $c\in[0,1]$, along with a \keyline{provenance} describing where this bound came from and how
\end{itemize}
}
\vspace{0.32cm}
\onslide<3->{%
\begin{ansonbox}{Endorsed gambles\hfill{\normalfont\color{Dust}\footnotesize Desirable gambles: Walley 1991}}
Equivalently: endorsing $E$ = accepting a \keyline{one-sided desirable gamble}
\par\smallskip
\centering
$\mu(\varphi\wedge\psi)\ge c\cdot\mu(\varphi)
\quad\Longleftrightarrow\quad
\mu(g_E)\ge 0,
\;\;
g_E=\mathbf 1_{\varphi\wedge\psi}-c\cdot\mathbf 1_{\varphi}$
\par\smallskip
{\small\color{Dust}pay $c$ for a ticket worth $1$ if $\psi$, bet is called off unless $\varphi$ obtains}
\par
\end{ansonbox}
}
\vspace{0.02cm}
\onslide<4->{%
{\fontsize{13.2}{16.8}\selectfont
The \keyline{endorsement book} $\mathcal E_t$ is the set of currently active endorsements\par}
}
\end{frame}

\begin{frame}[t]{Three ways to endorse a warrant}
\vspace{0.35cm}
{\fontsize{13.2}{16.8}\selectfont
Fix a $\rho$ and $n$ and abbreviate:\quad
$X=\alpha(n)$\qquad $A=\mathsf{App}_\rho(n)$\qquad $Y=\beta(n)$
\par}
\vspace{0.55cm}
\begin{ansonbox}{Three kinds of endorsements}
\vspace*{-1.5em}{\makebox[\dimexpr\linewidth-6.5pt\relax][r]{\color{Dust}\footnotesize\alt<2->{division-free form}{conditional form}}\par}\vspace{0.1em}
\smallskip
\centering
\begin{tabular}{@{}l@{\quad}l@{\qquad}c@{}}
$(E_P)$ & premise & $\mu(X)\ge c_P$\\[0.8em]
$(E_A)$ & applicability & \alt<2->{$\mu(X\wedge A)\ge c_A\cdot\mu(X)$}{$\mu(A\mid X)\ge c_A$}\\[0.8em]
$(E_S)$ & strength & \alt<2->{$\mu(X\wedge A\wedge Y)\ge c_S\cdot\mu(X\wedge A)$}{$\mu(Y\mid X\wedge A)\ge c_S$}
\end{tabular}
\par
\end{ansonbox}
\end{frame}

\begin{frame}[t]{Leverage propagation}
\vspace{-0.5cm}
\begin{columns}[T,onlytextwidth]
\begin{column}{0.54\textwidth}
\begin{plainbox}
For every compatible credence state $\mu$:
\par
{\fontsize{13}{14.2}\selectfont
\setlength\abovedisplayskip{1pt}\setlength\belowdisplayskip{1pt}
\[
\begin{aligned}
\mu(Y)\;&\ge\; \mu(X\wedge A\wedge Y) &&{\color{Dust}\text{\small (monotonicity)}}\\[0.12em]
\onslide<2->{&\ge\; c_S\,\mu(X\wedge A) &&{\color{Dust}(E_S)}}\\[0.12em]
\onslide<3->{&\ge\; c_S\,c_A\,\mu(X) &&{\color{Dust}(E_A)}}\\[0.12em]
\onslide<4->{&\ge\; c_S\,c_A\,c_P &&{\color{Dust}(E_P)}}
\end{aligned}
\]
}
\end{plainbox}
\end{column}
\begin{column}{0.38\textwidth}
\begin{outlineplain}
Suppose for some $\rho$ and $n$ we have all three endorsements:
\par\vspace{2pt}
\centering
$\mu(X)\ge c_P$\par\vspace{2pt}
$\mu(X\wedge A)\ge c_A\,\mu(X)$\par\vspace{2pt}
$\mu(X\wedge A\wedge Y)\ge c_S\,\mu(X\wedge A)$
\par
\end{outlineplain}
\end{column}
\end{columns}
\vspace{0.32cm}
\onslide<5->{%
{\centering\fontsize{13.2}{16.2}\selectfont
Endorsements of warrants multiplicatively propagate numerical bounds from one sentence to another
\par}
}
\end{frame}

\begin{frame}[t]{Credal regions}
\vspace{0.02cm}
{\fontsize{13.2}{16.8}\selectfont
\setlength\abovedisplayskip{4pt}\setlength\belowdisplayskip{2pt}
What happens when endorsements act jointly?
\onslide<2->{
\[
\mathcal C_t=\{\mu\in\Delta(\Omega_t):\ \mu(g_E)\ge0\ \text{for all }E\in\mathcal E_t\}
\]}
\par}
\vspace{-10pt}
\onslide<2->{%
{\fontsize{13.2}{16.8}\selectfont
``Every compatible credence state $\mu$'' means every $\mu\in\mathcal C_t$, distributions inside the \keyline{credal region} generated by the current endorsement book $\mathcal E_t$
\par}
}
\vspace{0.12cm}
\only<3>{%
\begin{ansonbox}{Three levels of demandingness}
\centering
\begin{minipage}[t]{0.31\linewidth}\centering\small
$\mathcal C_t=\Delta(\Omega_t)$\par\vspace{1.5pt}
the book forces nothing
\end{minipage}\hfill
{\setlength\fboxsep{2.5pt}\setlength\fboxrule{0.8pt}\color{SoftChalk}%
\fbox{\begin{minipage}[t]{\dimexpr0.34\linewidth-2\fboxsep-2\fboxrule\relax}\centering\small\color{Chalk}
$\emptyset\ne\mathcal C_t\subsetneq\Delta(\Omega_t)$\par\vspace{1.5pt}
constrained without being pinned
\end{minipage}}}\hfill
\begin{minipage}[t]{0.28\linewidth}\centering\small
$\mathcal C_t=\emptyset$\par\vspace{1.5pt}
the book demands irrationality\par\vspace{1.5pt}
{\color{Dust}``sure loss''}
\end{minipage}
\par
\end{ansonbox}
}%
\only<4->{%
\vspace{0.12cm}
\begin{outlineplain}
As the intersection of finitely many halfspaces, $\mathcal C_t$ is a polytope, which lets us use the math of \keyline{linear programming}
\end{outlineplain}
}
\end{frame}

\begin{frame}[t]{Credal intervals}
\vspace{-0.02cm}
{\fontsize{13.2}{16.8}\selectfont
Let's \textbf{\color{SoftChalk}project} onto a single sentence $\varphi$:
\par}
\vspace{0.12cm}
\onslide<2->{%
\begin{plainbox}
Lower and upper probabilities:
{\small\setlength\abovedisplayskip{3pt}\setlength\belowdisplayskip{2pt}
\[
\underline P_t(\varphi)=\min_{\mu\in\mathcal C_t}\mu(\varphi),
\qquad
\overline P_t(\varphi)=\max_{\mu\in\mathcal C_t}\mu(\varphi)
\]}
\vspace{0.55em}
which gives a \keyline{credal interval}
{\small\setlength\abovedisplayskip{3pt}\setlength\belowdisplayskip{0pt}
\[
\mu(\varphi)\in[\underline P_t(\varphi),\ \overline P_t(\varphi)]
\]}
\vspace{-0.6em}
\par
\end{plainbox}
}
\vspace{0.04cm}
\onslide<3->{%
{\fontsize{13.2}{16.8}\selectfont
Each endpoint is the value of a finite linear program over $\mathcal C_t$, so the min and max are attained
\par}
}
\end{frame}

\begin{frame}[t]{Worked example}
\vspace{-0.37cm}
{\fontsize{11.2}{13.2}\selectfont
At John's murder trial, he's accused of being seen in a bloody shirt:\ $\mathsf{Bloody}(\mathrm{John})\rightsquigarrow\mathsf{Guilty}(\mathrm{John})$
\par}
{\fontsize{11.0}{12.8}\selectfont
\setlength\topsep{0pt}\setlength\partopsep{0pt}
\begin{itemize}\setlength\itemsep{0.08em}
  \item $(c_P)$ ``Was John actually bloody?''
  \item $(c_A)$ ``Given that he was bloody, is that fact relevant to his guilt?''
  \item $(c_S)$ ``Given that he was bloody without extenuating circumstances, how strong a reason is that for his guilt?''
\end{itemize}
}
{\fontsize{12.2}{14.4}\selectfont
Take $(c_P,c_A,c_S)=(0.9,\,0.9,\,0.9)$:
\par}
{\setlength\abovedisplayskip{0pt}\setlength\belowdisplayskip{0pt}
\[
{\small
\mu(\mathsf{Guilty}(\mathrm{John}))\in[0,1]
\;\overset{\text{\small endorsements}}{\Longrightarrow}\;
\mu(\mathsf{Guilty}(\mathrm{John}))\in[\,0.729,\ 1\,]
}
\]}
\onslide<2->{%
\begin{ansonbox}{Remarks}
{\footnotesize
\begin{minipage}[t]{0.34\linewidth}
\textbullet\ LP says $\mu=0.729$ is attained
\end{minipage}\hfill
\begin{minipage}[t]{0.62\linewidth}
\textbullet\ A warrant like $X\rightsquigarrow\lnot Y$ can give us an upper bound below $1$
\end{minipage}\par}
\end{ansonbox}
}
\end{frame}

\begin{frame}[t]{LP duality}
\vspace{-0.3cm}
\only<1-3>{%
{\fontsize{12.8}{15.5}\selectfont
Linear programming provides a \keyline{dual certificate} $(\lambda,h)$ for $\underline P_t(\varphi)$:
\par}
{\setlength\abovedisplayskip{3pt}\setlength\belowdisplayskip{3pt}
\[
{\small
\mathbf 1_\varphi-{\color{ChalkPeach}c}\;=\;{\color{ChalkLilac}h}+\sum_{E\in\mathcal E_t}{\color{ChalkRose}\lambda_E}\,{\color{ChalkBlue}g_E},
\qquad {\color{ChalkRose}\lambda_E}\ge0,\ {\color{ChalkLilac}h}\ge0
}
\]}
{\footnotesize
{\color{ChalkPeach}$c$: the floor being certified}\ \ \ {\color{ChalkBlue}$g_E$: accepted gambles}\ \ \ {\color{ChalkRose}$\lambda_E$: weights}\ \ \ {\color{ChalkLilac}$h$: slack}
\par}
{\small\centering ``In every world, $\varphi$ pays out exactly enough to cover $c$ in cash plus a bundle of bets the book already accepts, with $h$ left over.''\par}
}%
\only<2>{%
\vspace{0.3cm}
\begin{outlineplain}
Apply $\mu\in\mathcal C_t$ to both sides:
{\setlength\abovedisplayskip{3pt}\setlength\belowdisplayskip{3pt}
\[
{\small
\mu(\varphi)-{\color{ChalkPeach}c}\;=\;\mu({\color{ChalkLilac}h})+\sum_{E\in\mathcal E_t}{\color{ChalkRose}\lambda_E}\,\mu({\color{ChalkBlue}g_E})\;\ge\;0
}
\]}
A certificate demonstrates $\mu(\varphi)\ge c$ for every compatible $\mu$
\end{outlineplain}
}%
\only<3>{%
\vspace{0.3cm}
\begin{ansonbox}{A place for whole-process feedback}
The $E$ with $\lambda_E>0$ are the reasons cited; the dual solution is the written opinion
\par\smallskip
$\lambda_E$ lets us receive feedback about mixtures of lots of different reasons
\end{ansonbox}
}%
\only<4->{%
\vspace{1.1cm}
\begin{ansonbox}{Off-the-shelf machinery\hfill{\normalfont\color{Dust}\small Lukasiewicz 2001}}
\smallskip
{\fontsize{12.5}{15.5}\selectfont
\begin{itemize}\setlength\itemsep{0.5em}
\item Our division-free bounds are \keyline{conditional constraints} $\mu(\psi\wedge\varphi)\ge\ell\cdot\mu(\varphi)$
\begin{itemize}
\item {\small An endorsement compiles to an already-understood object}
\end{itemize}
\item Two LPs compute $\underline P_t,\ \overline P_t$\qquad $\mathcal C_t\ne\emptyset\iff$ feasibility (Farkas' Lemma)
\begin{itemize}
\item {\small The natural problems come pre-solved}
\end{itemize}
\end{itemize}
}
\end{ansonbox}
}
\end{frame}

\begin{frame}[t]{Statics recap}
\only<1-4>{%
\vspace{-0.4cm}
{\setlength\abovedisplayskip{4pt}\setlength\belowdisplayskip{2pt}
\[
{\small
\underbrace{\rho:\ \alpha(n)\rightsquigarrow\beta(n)}_{\color{Dust}\text{relevance}}
\mskip\pipegapL\text{and}\mskip\pipegapR
\underbrace{\mathcal E_t}_{\color{Dust}\text{strength}}
\;\longrightarrow\;
\underbrace{\mathcal C_t}_{\color{Dust}\text{geometry}}
\;\longrightarrow\;
\underbrace{[\,\underline P_t(\varphi),\ \overline P_t(\varphi)\,]}_{\color{Dust}\text{leverage}}
\mskip\withgapL\text{with}\mskip\withgapR
\underbrace{(\lambda,h)}_{\color{Dust}\text{justification}}
}
\]}
\onslide<2->{%
{\fontsize{13.2}{16.8}\selectfont
Put in endorsed warrants, get out a justified credal interval
\par}
}
\vspace{0.01cm}
\onslide<3->{%
{\fontsize{11.5}{13.8}\selectfont
\setlength\topsep{1pt}
\begin{itemize}\setlength\itemsep{0.25em}
\item \textbf{\color{SoftChalk}``Middle way'' constraints:} real numerical constraints, not from settlement
\item \textbf{\color{SoftChalk}Reasoned disagreement:} you can accept the whole book and still disagree
\end{itemize}
}
}
\vspace{0.22cm}
\onslide<4->{%
\begin{ansonbox}{Epistemic graphs\hfill{\normalfont\color{Dust}\small Hunter--Polberg--Thimm 2020}}
{\fontsize{11}{13.5}\selectfont
\setlength\topsep{1pt}
\begin{itemize}\setlength\itemsep{0.15em}
\item Linear constraints on credences in probabilistic argumentation
\item Change over time only as one-shot belief revision: no settlement, no learning
\end{itemize}
}
\end{ansonbox}
}
}%
\only<5>{%
\vspace{1.5cm}
\begin{center}
{\fontsize{16}{22}\selectfont
Can we now move towards a \keyline{dynamics} of leverage?
\par}
\end{center}
}%
\end{frame}

\begin{frame}[t]{Provenance matters}
\vspace{-0.32cm}
{\fontsize{12}{14.2}\selectfont
Remember that every endorsement comes with a \keyline{provenance}, a ``paper trail'' identifying where the bound is coming from and what licensed it
\par}
\vspace{0.09cm}
\onslide<2->{%
{\fontsize{12}{14.2}\selectfont
\textbf{\color{SoftChalk}Tempting but bad idea:} discard the provenance and just work with the quantitative credal geometry
\par}
}
\vspace{0.28cm}
\onslide<3->{%
\begin{outlineplain}
{\footnotesize Two warrants $\rho:\alpha(n)\rightsquigarrow\beta(n)$ and $\sigma:\gamma(n)\rightsquigarrow\beta(n)$, each endorsed at the same $(c_P,\,c_A,\,c_S)$:\par}
\vspace{-9pt}
{\setlength\abovedisplayskip{0pt}\setlength\belowdisplayskip{0pt}
\[
{\footnotesize
\mu(\beta(n))\ge c_P\,c_A\,c_S\ \ {\color{Dust}\text{(via $\rho$)}}
\qquad\qquad
\mu(\beta(n))\ge c_P\,c_A\,c_S\ \ {\color{Dust}\text{(via $\sigma$)}}
}
\]}
\vspace{-1.15em}
\end{outlineplain}
}
\vspace{0.33cm}
\only<4-5>{%
\begin{minipage}[t]{0.475\linewidth}
\onslide<4->{%
\begin{ansonbox}{Conflict}
{\footnotesize Undercutter: $\delta(n)\rightsquigarrow\lnot\mathsf{App}_\rho(n)$\par\vspace{1.5pt}
The bound relaxes for $\rho$ and is unchanged for $\sigma$}
\end{ansonbox}
}
\end{minipage}\hfill
\begin{minipage}[t]{0.475\linewidth}
\onslide<5>{%
\begin{ansonbox}{Accrual}
{\footnotesize Should $\rho$ and $\sigma$ together mean a tighter bound than either alone?}
\end{ansonbox}
}
\end{minipage}
}%
\only<6->{%
\par\vspace*{0.55cm}
{\fontsize{12.8}{15.5}\selectfont
\textbf{\color{SoftChalk}Warrants with the same static geometry can have different dynamics}
\par}
}
\end{frame}

\begin{frame}[t]{Logical induction for normativity?}
\vspace{0.3cm}
{\fontsize{13.2}{16.8}\selectfont
Let's connect back to our motivation:
\par}
\vspace{0.9cm}
\onslide<2->{%
\begin{outlineplain}
{\fontsize{13.5}{17}\selectfont
A sharpened \keyline{desideratum}: any LI-based solution to Learning Normativity must have prices that respond to \textbf{\color{SoftChalk}leverage}
\par}
\end{outlineplain}
}
\end{frame}

\begin{frame}[t]{Warning}
\begin{center}
\vspace*{-0.10cm}
\includegraphics[height=2.35cm]{warning_symbol_yellow_transparent.png}\par
\vspace{0.60cm}
{\fontsize{17.5}{21}\selectfont Claims from this point on rest on LLM-assisted work}\par
\end{center}
\end{frame}

\begin{frame}[t]{Not so easy}
\vspace{-0.26cm}
{\fontsize{12.2}{14.8}\selectfont
We want the price to respect the book:\ \ $p_t(\mathsf{Guilty}(\mathrm{John}))\ge 0.729$
\par}
\vspace{0.06cm}
{\fontsize{11.2}{13.5}\selectfont
\setlength\topsep{1pt}
\begin{itemize}\setlength\itemsep{0.24em}
\item<2-> Doesn't vanilla LI already do this? \textbf{\color{SoftChalk}Not in general: only settlement and logic force prices}
\item<3-> Declare the sentence settled? \textbf{\color{SoftChalk}Goes against the whole point}
\item<4-> Add a trader representing the bound? \textbf{\color{SoftChalk}They'll go broke}
\item<5-> Make the trader rich? \textbf{\color{SoftChalk}The market can't tell wealth from evidence}
\end{itemize}
}
\vspace{0.10cm}
\onslide<6->{%
\begin{ansonbox}{Where can force enter?\quad{\normalfont\fontsize{10.8}{13.2}\selectfont Three LI objects we could modify:}}
{\fontsize{10.8}{13.2}\selectfont
\begin{minipage}[t]{0.30\linewidth}\centering
\textbf{\color{SoftChalk}Settlement}\\[0.25em]
Abandons the phenomenon
\end{minipage}\hfill
\begin{minipage}[t]{0.28\linewidth}\centering
\textbf{\color{SoftChalk}Traders}\\[0.25em]
Needs unlimited external funding
\end{minipage}\hfill
{\setlength{\fboxsep}{5pt}\setlength{\fboxrule}{1pt}%
\fcolorbox{SoftChalk}{BoardDeep}{\begin{minipage}[t]{0.26\linewidth}\centering
\textbf{\color{SoftChalk}Market Maker}\\[0.25em]
Most promising idea
\end{minipage}}}\par
\vspace{0.10em}}
\end{ansonbox}
\vspace{0.10cm}
}
\end{frame}

\begin{frame}[t]{Modifying the market maker}
\vspace{-0.16cm}
{\fontsize{12.2}{14.8}\selectfont
\textbf{\color{SoftChalk}New rule:}\ \ the market maker's quotes must respect the current book, $p_t\in\mathcal C_t$
\par}
\vspace{0.20cm}
{\fontsize{11.8}{14.4}\selectfont
Can traders exploit this modification?
\par}
\vspace{0.26cm}
\onslide<2->{%
% Display matches DECK_FACTS.md fact 2 / OF1_INDEPENDENT_AUDIT.md (packet, Aug 3).
% Deck-local renames for obviousness, to propagate at consolidation: B -> R (risk budget), D_N -> M_N (movement meter).
\begin{ansonbox}{Operative-Force Wealth Bound}
\vspace{0.10em}
{\centering
{\fontsize{13}{16}\selectfont
{\color{ChalkBlue}\textbf{winnings}}\;$\le$\;{\color{ChalkLilac}\textbf{slack}}\;$+$\;{\color{ChalkPeach}\textbf{rate}}\;$\times$\;({\color{SoftChalk}\textbf{risk}}\;$+$\;{\color{ChalkRose}\textbf{movement}})
\par}
\vspace{0.30em}
{\fontsize{9.5}{11.5}\selectfont
${\color{ChalkBlue}G_N(w)}\;\le\;{\color{ChalkLilac}\varepsilon_{1:N}}/{\color{ChalkPeach}\theta}\;+\;{\color{ChalkPeach}\bigl((1-\theta)/\theta\bigr)}\,\bigl({\color{SoftChalk}R}+{\color{ChalkRose}M_N}\bigr)$
\par}
\vspace{0.40em}
{\fontsize{11.2}{13.8}\selectfont
If the book always leaves the market maker at least a $\theta$-share of its freedom to quote, then it's not exploitable.
\par}\vspace{0.10em}}
\end{ansonbox}
}
\end{frame}

\begin{frame}[t]{Answerability}
\vspace{-0.06cm}
{\fontsize{11.6}{14.2}\selectfont
Now endorsements constrain prices, what constrains endorsements?
\par}
\vspace{0.20cm}
\onslide<2->{%
\begin{ansonbox}{Mechanism}
\vspace{0.10em}
{\fontsize{10.8}{13.4}\selectfont
Challengers post a bond and issue \keyline{challenges} against particular commitments of the book, governed by checkable rules. The book pays \keyline{charges} to the challenger each period until it \keyline{answers}: refuting the challenge under the rules, or revising through an authorized route. A book is \keyline{answerable} if every challenge eventually ends, every revision travels an authorized route, and the total charges ever paid are finite.
\par}\vspace{0.10em}
\end{ansonbox}
}
\vspace{0.26cm}
\onslide<3->{%
{\centering\fontsize{11.8}{14.4}\selectfont
\textbf{\color{SoftChalk}Diachronic criterion:} you must be answerable to your past self
\par}
}
\end{frame}

\begin{frame}[t]{Learning without ground truth}
\vspace{-0.06cm}
{\fontsize{11.6}{14.2}\selectfont
What constitutes improvement when there is no objective standard?
\par}
\vspace{0.20cm}
\onslide<2->{%
\begin{ansonbox}{$\Phi$-regret\hfill{\normalfont\color{Dust}\small Greenwald--Jafari 2003}}
\vspace{0.02em}
{\fontsize{10.8}{13.4}\selectfont
How much loss $\ell_t$ can be reduced from action $a_t$ within a single edit class $\Phi$:
\par}
{\setlength\abovedisplayskip{1pt}\setlength\belowdisplayskip{0pt}
\[
{\small
R_T^{\Phi}\;=\;\sup_{\varphi\in\Phi}\ \sum_{t\le T}\bigl[\,\ell_t(a_t)-\ell_t(\varphi(a_t))\,\bigr]
}
\]}
\end{ansonbox}
}
\vspace{0.20cm}
\onslide<3->{%
{\fontsize{10.9}{13.4}\selectfont
Take $\ell_t$ = charges owed and $\Phi$ = the \keyline{lawful} edits of the book
\par}
\vspace{0.20cm}
{\fontsize{10.9}{13.4}\selectfont
\textbf{\color{SoftChalk}Lawfulness:} part global condition (\keyline{answerability}) and part local condition (\keyline{reasons-responsiveness})
\par}
}
\end{frame}

\begin{frame}[t]{Status of LLM work}
\vspace{0.06cm}
{\fontsize{10.6}{13.2}\selectfont
\setlength\topsep{1pt}
\begin{itemize}\setlength\itemsep{0.34em}
\item \textbf{\color{SoftChalk}Force:} not exploitable while a core of quoting freedom is preserved
\item \textbf{\color{SoftChalk}Answerability:} an explicit policy is answerable from any authorized start; minimum funding exact
\item \textbf{\color{SoftChalk}Composition:} both pressures on one history; no interference as long as funding is kept separate
\item \textbf{\color{SoftChalk}Identity:} obligations survive vocabulary change; nothing merges silently
\item \textbf{\color{SoftChalk}Demand:} bounded liability forces rulings to track admitted arrivals
\item \textbf{\color{SoftChalk}Settlement:} an interface world-reports must meet before they get force
\item \textbf{\color{SoftChalk}Learning:} candidate criterion; staged proof plan
\end{itemize}
}
\vspace{0.18cm}
\onslide<2->{%
\begin{ansonbox}{In the repo}
\vspace{0.06em}
{\fontsize{9.8}{12.2}\selectfont
Finite models in Python; hand proofs of theorems; one script checks every claim
\par}
\end{ansonbox}
}
\end{frame}

\begin{frame}[t]{Open questions}
\vspace{0.12cm}
{\fontsize{12.2}{15.4}\selectfont
\setlength\topsep{1pt}
\begin{itemize}\setlength\itemsep{1.25em}
\item What precisely is ``reasons-responsiveness''?
\item What's the real relationship to LI? Is this a new LI-like construction, or an extension of LI?
\item What should demand optimize for? The existing \keyline{docket} story just gives a feasible region
\item Where exactly do endorsements come from? Candidate answer: \keyline{inquiry}
\end{itemize}
}
\end{frame}

\begin{frame}[t]{Conclusion}
\vspace{1.55cm}
{\setlength{\fboxsep}{6pt}\setlength{\fboxrule}{1pt}%
\centering
\onslide<1->{\fcolorbox{SoftChalk}{Board}{\begin{minipage}[c]{\dimexpr0.26\linewidth-2\fboxsep-2\fboxrule\relax}
{\fontsize{10.5}{13.2}\selectfont\color{Chalk}\raggedright\hyphenpenalty=10000\exhyphenpenalty=10000
The \textbf{\color{SoftChalk}concept} of leverage is what's missing from logical induction for normativity\par}
\end{minipage}}}%
\onslide<2->{\ {\Large$\longrightarrow$}\ }%
\onslide<2->{\fcolorbox{SoftChalk}{Board}{\begin{minipage}[c]{\dimexpr0.26\linewidth-2\fboxsep-2\fboxrule\relax}
{\fontsize{10.5}{13.2}\selectfont\color{Chalk}\raggedright\hyphenpenalty=10000\exhyphenpenalty=10000
The \textbf{\color{SoftChalk}statics} of leverage use known math for a concrete desideratum\par}
\end{minipage}}}%
\onslide<3->{\ {\Large$\longrightarrow$}\ }%
\onslide<3->{\fcolorbox{SoftChalk}{Board}{\begin{minipage}[c]{\dimexpr0.26\linewidth-2\fboxsep-2\fboxrule\relax}
{\fontsize{10.5}{13.2}\selectfont\color{Chalk}\raggedright\hyphenpenalty=10000\exhyphenpenalty=10000
The \textbf{\color{SoftChalk}dynamics} of leverage give rise to hard, but tractable problems with open design questions\par}
\end{minipage}}}\par}
\end{frame}

\begin{frame}[t]{References}
\vspace{0.15cm}
{\small
\begin{itemize}\setlength\itemsep{0.45em}
  \item A. Demski. \emph{Learning Normativity: A Research Agenda.} Alignment Forum, 2020.
  \item A. Greenwald and A. Jafari. A general class of no-regret learning algorithms and game-theoretic equilibria. \emph{COLT}, 2003.
  \item A. Hunter, S. Polberg, and M. Thimm. Epistemic graphs for representing and reasoning with positive and negative influences of arguments. \emph{Artificial Intelligence} 281:103236, 2020.
  \item T. Lukasiewicz. Probabilistic logic programming with conditional constraints. \emph{ACM Transactions on Computational Logic} 2(3):289--339, 2001.
  \item S. Nair and J. Horty. The logic of reasons. In D. Star (ed.), \emph{The Oxford Handbook of Reasons and Normativity}. Oxford University Press, 2018.
  \item P. Walley. \emph{Statistical Reasoning with Imprecise Probabilities.} Chapman \& Hall, 1991.
\end{itemize}
}
\end{frame}

\end{document}